% Generated draft: requires mathematical and source review.
\documentclass[11pt]{article}
\usepackage[margin=25mm]{geometry}
\usepackage{fontspec}
\setmainfont{Latin Modern Roman}
\usepackage{amsmath,amssymb,mathtools,unicode-math}
\setmathfont{Latin Modern Math}
\usepackage{enumitem,needspace,xurl,hyperref}
\setlength{\parindent}{0pt}
\setlength{\parskip}{5pt}
\setlength{\emergencystretch}{3em}
\newcommand{\N}{\mathbb{N}}
\newcommand{\Z}{\mathbb{Z}}
\newcommand{\Q}{\mathbb{Q}}
\newcommand{\R}{\mathbb{R}}
\newcommand{\C}{\mathbb{C}}
\newcommand{\F}{\mathbb{F}}
\newcommand{\A}{\mathbb{A}}
\newcommand{\PP}{\mathbb{P}}
\newcommand{\E}{\mathbb{E}}
\newcommand{\eps}{\varepsilon}
\newcommand{\dd}{\,\mathrm{d}}
\begin{document}
\section*{UnsolvedMath ID: 800954}
\textbf{Source problem code:} catalogue entry
\subsection{Definitions and mathematical statement}
The supplied excerpt does not specify the Riemannian manifold, completeness assumptions, invariant measure, or hypotheses ensuring ergodicity. Consequently, the following gives a conditional interpretation of its notation, with the unresolved qualifications stated explicitly.

Let \((M,g)\) be a connected Riemannian manifold and let \(\pi:\widetilde M\to M\) be its universal covering, equipped with the lifted Riemannian metric. Write
\[
SM=\{v\in TM:g(v,v)=1\}
\]
for the unit tangent bundle. If the metric is geodesically complete, let \(\gamma_v:\R\to M\) be the geodesic with initial velocity \(v\), and let \(\widetilde\gamma_v:\R\to\widetilde M\) be a lift. Completeness is an additional interpretive assumption here, not a hypothesis supplied by the excerpt.

In this setting the first two displayed sets in the source are
\begin{align*}
T_b&=\{v\in SM: \text{every lift of }\gamma_v\text{ has bounded image in }\widetilde M\},\\
T_p&=\{v\in SM: \text{every lift of }\gamma_v\text{ has unbounded image and approaches infinity properly}\}.
\end{align*}
Here boundedness refers to the lifted Riemannian distance. One standard interpretation of ``approaches infinity properly'' is that each lift is a proper map: for every compact \(K\subseteq\widetilde M\),
\[
\widetilde\gamma_v^{-1}(K)\text{ is compact in }\R.
\]
For a geodesic defined on all of \(\R\), this requires eventual avoidance of each compact set in both time directions. The excerpt does not define the phrase or specify whether both time directions are intended, so this interpretation requires confirmation.

The source defines the third set literally by
\[
T_i=TM\setminus(T_b\cup T_p).
\]
This uses \(TM\), although the first two sets lie in \(SM\). In particular, the displayed definition includes every nonunit tangent vector in \(T_i\). If the intended phase space is \(SM\), the corresponding third class would instead be
\[
T_i^{SM}=SM\setminus(T_b\cup T_p).
\]
These are distinct definitions, and the supplied record does not resolve which is intended.

\textbf{Source question.} The excerpt asserts that the three sets are measurable and invariant under the geodesic flow and that ergodicity implies that one has full measure. It asks: which of these sets has full measure?

More precisely, conditional on the intended phase space being \(SM\), the preceding interpretation of \(T_p\), and a specified geodesic-flow-invariant ergodic probability measure \(\mu\) for which these sets are measurable, the question is which member of
\[
\{T_b,T_p,T_i^{SM}\}
\]
has \(\mu\)-measure \(1\). The supplied excerpt does not identify \(M\), \(g\), or \(\mu\), or supply the assumptions that would make this conditional formulation the intended question. Its measurability and ergodicity assertions are therefore retained as source assertions requiring their missing context.
\subsection{Short English statement}
The question asks whether almost every unit-speed geodesic has bounded lifts to the universal cover, lifts that escape properly, or lifts in the remaining class. The manifold, measure, ergodicity hypotheses, and meaning of proper escape are missing. The source also defines the remaining class in \(TM\), although its first two classes lie in \(SM\); this discrepancy requires review.
\subsection{Sources}
\begin{itemize}
\item Supplied local UnsolvedMath registry record, permanent ID \(800954\), exact source code \texttt{catalogue entry}; upstream revision \texttt{d371ee917ac9647fbb06cd6fc529deec6f3dda82}. The registry labels the problem \emph{open}; this is supplied metadata, not an independently verified status assessment.
\item Burns and Matveev, \emph{Open Problems and Questions about Geodesics}, Question 34, listed in the supplied record as 2019. Supplied source URL: \url{https://arxiv.org/abs/1308.5417}. This citation and its connection to the excerpt have not been independently checked.
\item The supplied title and background attribute the question or its context to Benjamin Schmidt. The registry's \texttt{proposed\_by} field lists Burns and Matveev; both forms of attribution are retained without resolving their respective roles.
\item The supplied provenance reports access on 2026-10-16 and extraction by \texttt{arxiv\_latex}. This is a supplied access record, not a claim of access or verification for this draft.
\end{itemize}
\end{document}
